\section{Related work}
\label{sec:related}

\paragraph{Benchmarks that hand the agent its data.}
The dominant shape gives an agent a repository or a dataset and scores the
artifact it produces. SWE-bench~\citep{jimenez2024swebench} scores patches
against held-out tests; AgentBench~\citep{liu2024agentbench} covers interactive
environments more broadly. Closest to our setting are the ML-engineering
suites: MLE-bench~\citep{chan2024mlebench} runs agents on Kaggle competitions
and grades against the private leaderboard;
MLAgentBench~\citep{huang2024mlagentbench} scores iterative experimentation on
fixed task datasets; DSBench~\citep{jing2024dsbench} draws on data-science
competitions.

A second line targets science specifically.
ScienceAgentBench~\citep{chen2025scienceagentbench} assembles data-driven
discovery tasks from published papers;
DiscoveryBench~\citep{majumder2025discoverybench} formalises hypothesis search
over provided datasets; BixBench~\citep{mitchener2025bixbench} does the same for
computational biology; CORE-Bench~\citep{siegel2024corebench} scores
computational reproducibility; PaperBench~\citep{starace2025paperbench} scores
replication of ML papers; RE-Bench~\citep{wijk2024rebench} pits agents against
human experts on open-ended research engineering.

These differ enormously in domain and open-endedness, and in one respect they
agree: the training data exists in advance and is identical for every agent.
That is what makes them clean comparisons --- and it is also what removes
acquisition from the set of decisions being measured.

\paragraph{The literature the task asks an agent to rediscover.}
Deciding where to measure next is of course a mature field. Active
learning~\citep{settles2009active} and Bayesian experimental
design~\citep{chaloner1995bayesian} supply the strategies a competent agent
might reach for --- uncertainty sampling, residual-targeted acquisition,
variance reduction under a probabilistic model --- and space-filling designs
such as Latin hypercubes~\citep{mckay1979lhs} supply the baseline those
strategies have to beat. We contribute nothing to that literature. What we
measure is whether an agent reaches for any of it unprompted, implements what it
says it implemented, and spends a hard budget sensibly
(Section~\ref{sec:agentlogic}).

\paragraph{What we add, stated narrowly.}
The gap is not that nobody acquires data. It is that \emph{agentic coding
benchmarks do not make the acquisition policy part of what is scored}. To our
knowledge, as of mid-2026, no existing agent benchmark requires the agent to
(i) drive an expensive, failure-prone numerical solver for which it has not been
handed a working harness, (ii) allocate a hard budget across adaptive rounds,
and (iii) be scored on a frozen set it never sees, with a zero-LLM scripted
policy running inside the identical contract as the reference point.

Two recent works come closest. AstaBench~\citep{bragg2025astabench} assembles a
large scientific-agent suite and is motivated by the same complaint about
confounders --- model cost and tool access varying across compared systems ---
that motivates our model control; it supplies the tools and the data rather than
scoring acquisition. And a position paper argues that the execution harness,
not the model it wraps, is often the binding constraint on agent
performance~\citep{zhang2026harness}. Our campaign is a direct empirical test of
that claim: one model, four harnesses, and the spread measured
(Section~\ref{sec:harness-effect}).

\paragraph{The agent frameworks we evaluate.}
These are existing systems, not contributions of this paper:
URSA~\citep{ursa2025}, a planner/executor pair built on LangGraph, part of the
LangChain project~\citep{langchain2022}; DSPy~\citep{khattab2024dspy}, whose
typed modules we also use for the \Lone{} decision layer and whose GEPA
optimiser~\citep{agrawal2025gepa} we use for prompt optimisation; and two
commercial command-line coding agents. We evaluated them as we found them; none
of their authors was involved in or consulted about this evaluation.
